\section{Future Work}
Building on the promising results of this study, several research directions warrant further investigation to advance our understanding of LLM-based POI personalization and its practical deployment in travel platforms.

\textbf{Diverse User Populations.} Future studies should recruit participants with greater demographic diversity, particularly including older adults, non-technical professionals, and individuals from varied cultural backgrounds. Understanding how different age groups and cultural contexts influence receptiveness to AI-generated travel content would provide crucial insights for global platform deployment. Additionally, examining how travel expertise level moderates preferences for personalized versus generic descriptions could inform adaptive personalization strategies.

\textbf{Longitudinal Evaluation.} This study captured participant perceptions at a single point in time. Longitudinal research tracking how user preferences and trust in AI-generated descriptions evolve with repeated exposure would illuminate whether positive initial impressions persist or fade over time. Such studies could also investigate whether users develop preferences for specific personalization styles or become more discerning about AI-generated content quality with experience.

\textbf{Behavioral Outcomes.} While our study measured perceptions and stated preferences, future research should examine whether personalized descriptions influence actual behavior. Do personalized POI descriptions increase visit likelihood, booking rates, or satisfaction with chosen destinations? Field studies conducted in partnership with travel platforms could track conversion metrics and post-visit satisfaction to assess real-world impact.

\textbf{Personalization Granularity.} This study implemented personalization at a relatively coarse level based on stated interests and travel styles. Future work could explore finer-grained personalization considering factors such as trip purpose (business vs. leisure), travel companions (solo, couple, family), budget constraints, accessibility needs, and contextual factors like time of year or local events. Investigating optimal personalization granularity would help balance customization benefits against complexity and potential privacy concerns.

\textbf{Transparency and Explainability.} Our study did not inform participants which descriptions were AI-generated during the evaluation phase. Future research should examine how transparency about AI involvement affects user perceptions and preferences. Additionally, providing explanations for why specific aspects were emphasized in personalized descriptions might enhance trust and user understanding of the personalization process.

\textbf{Comparison with Recommendation Systems.} While this study focused on adapting descriptions of given POIs, future work could integrate personalized descriptions with POI recommendation systems. Investigating how personalized descriptions complement personalized recommendations could optimize the entire discovery-to-booking journey.

\textbf{Mitigation of Filter Bubbles.} Future research should examine potential negative effects of over-personalization, such as filter bubbles that limit serendipitous discovery of POIs outside user-stated interests. Developing strategies to balance personalization with diversity and novelty would ensure that adaptive systems enhance rather than constrain travel exploration.
